%==============================================================================
%  Chladni Plate Inverse-Design — academic paper (English)
%  Compile:  xelatex chladni_paper_en.tex  (run twice for references)
%            or  pdflatex chladni_paper_en.tex  (run twice)
%  The experimental section is intentionally left as a scaffold (to be completed).
%==============================================================================
\documentclass[11pt,a4paper]{article}

\usepackage[T1]{fontenc}
\usepackage{amsmath,amssymb,amsfonts}
\usepackage{bm}
\usepackage[a4paper,margin=2.4cm]{geometry}
\usepackage{booktabs}
\usepackage{array}
\usepackage{enumitem}
\usepackage{microtype}
\usepackage{caption}
\usepackage{xcolor}
\usepackage[hidelinks,breaklinks]{hyperref}
\usepackage{titlesec}

\titleformat{\section}{\normalfont\large\bfseries}{\thesection}{0.6em}{}
\titleformat{\subsection}{\normalfont\normalsize\bfseries}{\thesubsection}{0.6em}{}
\setlist{itemsep=2pt,topsep=3pt,parsep=0pt}

\newcommand{\Done}{D_1}     % placeholder guard, not used
\newcommand{\Aone}{A_1}
\newcommand{\Dfour}{D_4}
\newcommand{\matH}{\mathbf{H}}
\newcommand{\matK}{\mathbf{K}}
\newcommand{\matM}{\mathbf{M}}

\title{\bfseries Physical Reachability Limits of Pattern Inverse-Design for
Centrally-Driven Chladni Plates:\\[2pt]
\large A Symmetry-Constrained, Differentiable-Surrogate-to-FEA Framework with
Manufacturable Output}

\author{%
Chladni Plate Inverse-Design Team\\
\small Department of Physics, Imperial College London\\
\small Term~3 Summer Project (first-year undergraduate)\\
\small \texttt{ChladniPlateInverseSystem}}

\date{31 May 2026}

\begin{document}
\maketitle

%------------------------------------------------------------------ Abstract
\begin{abstract}
\noindent
We study a classic and severely constrained inverse problem: given a target
image, recover a plate (a spatial thickness field, and optionally a per-cell
fibre orientation) together with a set of drive frequencies whose Chladni sand
pattern approximates the target. The problem is posed under a fixed,
non-negotiable hardware contract: a $150\,\text{mm}$ square plate discretised by
a $15\times15$ thickness grid, an $8\,\text{mm}$ central clamp, free edges, and a
\emph{single} central point excitation with a frequency sweep, with a
high-fidelity finite-element model (COMSOL Multiphysics with
LiveLink\texttrademark{} for MATLAB) as the primary source of truth. Across an
extended development effort spanning three algorithmic ``eras'', roughly twenty
algorithm lines and several hundred finite-element candidates, we converge on
four central, and largely counter-intuitive, conclusions. (i)~Arbitrary
custom-pattern inversion is \emph{physically impossible} rather than an
engineering shortfall: the dihedral symmetry group $\Dfour$ of the configuration
forces a central point drive to couple only into the $\Aone$ irreducible
representation, so the fraction of target energy lying in $\Aone$ is a hard
reachability ceiling. (ii)~Material anisotropy (the fibre orientation $\theta$)
is essentially inert in the working band, because the inertial term dominates the
stiffness term by several orders of magnitude and $\theta$ enters the physics only
through stiffness. (iii)~The gap between a differentiable thin-plate surrogate and
the thick-plate finite-element model is the dominant practical ceiling, and is a
systematic physical difference rather than numerical noise. (iv)~The honest
deliverable is therefore not ``arbitrary inversion'' but a reproducible,
manufacturable \emph{system}: a differentiable surrogate used as a compass, a
two-phase finite-element production pipeline (modal selection followed by
trust-region refinement), and a directly 3D-printable export. We develop the
theory and the methodology in full; the quantitative experimental validation is
reported separately and the corresponding section here is left as a scaffold.

\medskip
\noindent\textbf{Keywords:} Chladni figures; inverse design; structural dynamics;
plate vibration; dihedral symmetry; irreducible representations; differentiable
surrogate; topology/thickness optimisation; finite-element validation;
additive manufacturing.
\end{abstract}

%------------------------------------------------------------------ 1
\section{Introduction}

\subsection{Background and motivation}
Chladni figures are the patterns traced by fine powder on a vibrating plate: the
powder migrates away from antinodes and accumulates along nodal lines where the
out-of-plane displacement is minimal, rendering an eigenmode visible to the naked
eye~\cite{chladni1787,waller1939}. The \emph{forward} problem---predicting the
nodal pattern from the plate geometry, material and drive frequency---is, in
principle, a standard exercise in linear structural dynamics. The
\emph{inverse} problem we consider here is far harder: \emph{given a desired
target image, recover a plate design and a drive frequency whose Chladni pattern
reproduces it.} Inverse design of this kind sits at the intersection of
structural-dynamics eigenvalue placement, shape/topology optimisation, and
manufacturability, and it is an attractive vehicle for studying how far a purely
computational pipeline can be pushed against a genuinely fixed piece of
hardware.

\subsection{Problem statement}
We fix a square plate excited at its centre by a single point source and ask
whether, by spatially varying the plate thickness (and, as a tempting extra
degree of freedom, the local fibre orientation of a 3D-printed part) and by
choosing one or more sweep frequencies, the resulting nodal pattern can be made to
resemble an arbitrary target such as a letter pair, a cross, a diagonal, or a
ring. The constraints (Section~\ref{sec:contract}) are immutable for the entire
study; every algorithm must be designed \emph{around} them rather than relaxing
them.

\subsection{Contributions and central thesis}
Our central thesis is that the most valuable output of such a study is not a
demonstration that ``any pattern can be produced''---which we show does not exist
physically---but a precise characterisation of \emph{where the reachability
boundary lies, why it lies there, and how to engineer the boundary-hugging
optimum into a deliverable product.} Concretely, this paper contributes:
\begin{enumerate}[label=(\arabic*)]
\item a \textbf{symmetry-based reachability theory} that identifies the $\Aone$
irreducible representation of the dihedral group $\Dfour$ as the only channel a
central point drive can excite, yielding a closed-form, target-specific upper
bound on what is achievable (Section~\ref{sec:theory});
\item a \textbf{scaling analysis of anisotropy} showing that the fibre
orientation $\theta$ is a parasitic degree of freedom in the working band, with a
mass--stiffness scale separation of several orders of magnitude that renders its
sensitivity negligible (Section~\ref{sec:aniso});
\item a \textbf{differentiable orthotropic-plate surrogate} with adjoint/autograd
sensitivities, framed explicitly as a ``compass, not a judge''
(Section~\ref{sec:surrogate});
\item a \textbf{recognisability metric} grounded in the actual powder-deposition
physics rather than pixel overlap (Section~\ref{sec:metric});
\item a \textbf{two-phase finite-element production pipeline} (modal selection
followed by trust-region refinement) that manages the surrogate--FEA gap
explicitly (Section~\ref{sec:pipeline}); and
\item a \textbf{manufacturable export and software system} that turns a design
into a printable part (Section~\ref{sec:software}).
\end{enumerate}
Throughout, we adopt a deliberately honest, project-report register: negative
results, retracted assumptions, and methodological mistakes are reported as
first-class findings, because in a fixed-hardware inverse problem they delimit the
feasible region as sharply as any positive result.

\subsection{Organisation}
Section~\ref{sec:related} reviews directly comparable prior art.
Section~\ref{sec:contract} formalises the problem and the hardware contract.
Section~\ref{sec:theory} develops the symmetry ceiling.
Section~\ref{sec:method} presents the full methodology (surrogate, anisotropy
analysis, metric, pipeline, export, software). Section~\ref{sec:experiments} is
the experimental section and is intentionally left as a scaffold to be completed.
Section~\ref{sec:discussion} discusses the lessons; Section~\ref{sec:limits}
states the limitations and future work; Section~\ref{sec:conclusion} concludes.

%------------------------------------------------------------------ 2
\section{Related work}\label{sec:related}

\paragraph{Topology optimisation of Chladni / modal patterns.}
The most directly comparable prior art is the 2024 study on topology-optimisation
design of Chladni patterns for vibration-mode manipulability~\cite{acta2024},
which uses density-based topology optimisation (SIMP) with an objective that
minimises the error between computed eigenvectors and a target Chladni pattern,
inversely designing material distributions for single- and multi-order patterns.
Its eigenvector-error objective is cognate to our own scoring evolution, and its
SIMP density variable is close in spirit to our density/thickness field. The key
difference is that it is \emph{not} bound by a single-central-excitation $\Dfour$
constraint---the very source of our $\Aone$ ceiling---so its reachable pattern
space is intrinsically larger. The methodological lineage of mode-targeting SIMP
traces back to topology optimisation of resonating structures~\cite{tcherniak2002}
and, more broadly, to eigenvalue topology
optimisation~\cite{pedersen2000,achtziger2007} within the classical
density-based framework~\cite{bendsoe1989,bendsoe1999,bendsoe2003}.

\paragraph{Energy-landscape formulations.}
An open-source line of work frames Chladni inverse design as an
``energy-landscape'' problem in which a nonnegative mixture of squared mode shapes,
$E(x,y)=\sum_k b_k\,\varphi_k(x,y)^2$, is shaped toward the
target~\cite{bellette}. This matches our amplitude-valley loss and independently
confirms a structural bias: nonnegative squared mixtures prefer boxes, bands,
grids, crosses and large connected glyphs, and resist fine, separated features---a
bias strongly consistent with what our symmetry analysis and powder physics
predict.

\paragraph{Cymatics and wave control in structured plates.}
Reshaping flexural-wave fields and nodal lines through plate microstructure has
been demonstrated in the context of flexural cloaking~\cite{misseroni2016}, and
forced-response Chladni patterns have been related to maximum-entropy states of
an inhomogeneous Helmholtz equation~\cite{tuanchen}, both of which inform the
distinction between forced response and free vibration that we revisit in
Section~\ref{sec:gap}.

\paragraph{Hardware-driven control.}
Finally, multi-actuator active control of particles on a Chladni
plate~\cite{zhou2016} demonstrates that genuinely breaking the symmetry ceiling
requires acting on the \emph{hardware}---additional shakers or off-centre
excitation---a conclusion that our theory reaches from the opposite direction and
that frames our recommendations for future work. Level-set shape
optimisation~\cite{allaire2004} is noted as an alternative geometry representation
that we did not adopt, preferring a density/thickness field.

%------------------------------------------------------------------ 3
\section{Problem formulation and hardware contract}\label{sec:contract}

\subsection{Forward model}
Let $\Omega=[0,L]^2$ be the mid-surface of a thin plate of spatially varying
thickness $h(\bm{x})$, clamped on a small central disc and free on all four
edges. Under the Kirchhoff--Love thin-plate assumption, the free-vibration
eigenproblem is the biharmonic-type generalised problem
\begin{equation}
\matK(\bm{h})\,\bm{\phi}_i \;=\; \omega_i^2\,\matM(\bm{h})\,\bm{\phi}_i,
\label{eq:eig}
\end{equation}
where $\matK$ and $\matM$ are the (discretised) stiffness and mass operators,
$\omega_i$ the angular eigenfrequencies and $\bm{\phi}_i$ the mode shapes. A
high-fidelity description retains transverse shear via the Mindlin--Reissner
thick-plate theory~\cite{reissner1945,mindlin1951}, which is the model used inside
the finite-element solver. Under a steady single-point harmonic excitation of
angular frequency $\omega$ and load vector $\bm{f}$, the forced response is
\begin{equation}
\bm{u}(\omega)=\big(\matK-\omega^2\matM + \mathrm{i}\,\omega\,\mathbf{C}\big)^{-1}\bm{f},
\label{eq:forced}
\end{equation}
with $\mathbf{C}$ a damping operator. The Chladni pattern is the locus of small
$|\bm{u}|$: powder is deposited where the out-of-plane amplitude is minimal,
which we model by a powder-density functional that concentrates mass on nodal
lines (Section~\ref{sec:metric}).

\subsection{Inverse problem}
The inverse problem seeks design variables that make the deposited pattern
resemble a target image $T$. The primary design variable is the thickness field,
discretised on the mandated grid,
\begin{equation}
\matH=\{h_{mn}\}_{m,n=1}^{15}\in\mathbb{R}^{15\times15},
\end{equation}
optionally augmented by a per-cell fibre orientation field $\bm{\theta}$ and by a
small set of $K$ drive frequencies $\{\omega_k\}$ with weights $\{w_k\}$. The
objective is a recognisability functional $\mathcal{R}(\,\cdot\,;T)$ defined in
Section~\ref{sec:metric}.

\subsection{Immutable hardware contract}
Table~\ref{tab:contract} lists the constraints that remained fixed throughout the
study and that define the physical ceiling of the problem.

\begin{table}[h]
\centering
\caption{Immutable hardware contract. Every algorithm is designed around these.}
\label{tab:contract}
\small
\begin{tabular}{@{}lll@{}}
\toprule
Item & Value & Source \\
\midrule
Design grid          & $15\times15$ thickness field        & FEA \texttt{H.csv} contract \\
Plate size           & $150\,\text{mm}\times150\,\text{mm}$ & one-shot print boundary \\
Central clamp        & radius $8\,\text{mm}$ disc          & standard rig \\
Edges                & free on all four sides              & free vibration \\
Excitation           & single central point + sweep        & no multi-shaker / added mass \\
Primary feedback     & FEA eigenfrequency / forced response& LiveLink with MATLAB \\
Thickness range      & $0.6$--$2.0\,\text{mm}$, neighbour $\Delta\le1.0\,\text{mm}$ & manufacturability \\
\bottomrule
\end{tabular}
\end{table}

%------------------------------------------------------------------ 4
\section{The symmetry ceiling}\label{sec:theory}

\subsection{Dihedral symmetry of the configuration}
A square plate with a central clamp and a central point drive is invariant under
the dihedral group $\Dfour$ of order eight (four rotations and four
reflections). Because both the operator pencil $(\matK,\matM)$ and the load
$\bm{f}$ of a \emph{centred} excitation are $\Dfour$-symmetric, the entire
steady-state response must transform according to the representation theory of
$\Dfour$~\cite{tinkham1964}. The group has five irreducible representations
(irreps): the one-dimensional $A_1,A_2,B_1,B_2$ and the two-dimensional $E$.

\subsection{Central excitation couples only to \texorpdfstring{$\Aone$}{A1}}
The load of a point force applied exactly at the (clamped) centre is itself
invariant under every element of $\Dfour$; it therefore lies entirely in the
trivial representation $\Aone$. By orthogonality of irreps, such a load can only
excite response components that also transform as $\Aone$: the projection of the
forced response onto any non-$\Aone$ subspace vanishes identically. We implement
this projection explicitly as a $\Dfour$ irreducible-representation
decomposition of any candidate field. The practical consequence is a clean,
\emph{closed-form} reachability statement:
\begin{equation}
\boxed{\;\text{reachability}(T)\;=\;\frac{\lVert \mathcal{P}_{\Aone} T\rVert^2}
{\lVert T\rVert^2}\;,}
\label{eq:reach}
\end{equation}
where $\mathcal{P}_{\Aone}$ projects the target onto the $\Aone$ subspace. Energy
of the target outside $\Aone$ cannot be reproduced \emph{no matter how strong the
algorithm}, because the central drive provides no channel into those irreps.

\subsection{Worked reachability of representative targets}
Equation~\eqref{eq:reach} is analytic and can be evaluated directly on each
target. Table~\ref{tab:reach} reports the resulting $\Aone$ fractions. An X-shape
lives entirely in $\Aone$ and is in principle fully reachable; a single diagonal
is only half-reachable (its complement lies in $B_2$); and the ``IC'' letter pair
places the majority of its energy in $E$ and $B_1$, so most of its signal is
\emph{structurally} inadmissible. This furnishes, for the first time in the
project, an exact mathematical answer to the question ``why does IC never appear''
that is independent of any particular optimiser.

\begin{table}[h]
\centering
\caption{Analytic reachability (Eq.~\eqref{eq:reach}): the $\Aone$ energy
fraction is a hard upper bound a central drive can hope to reproduce.}
\label{tab:reach}
\small
\begin{tabular}{@{}lccl@{}}
\toprule
Target & $\Aone$ fraction & Outside $\Aone$ (unreachable) & Meaning \\
\midrule
X-shape          & $100.0\%$ & $0\%$            & theoretically fully reachable \\
Single diagonal  & $51.7\%$  & $48.3\%$ ($B_2$) & half reachable \\
IC letters       & $41.9\%$  & $58.1\%$ ($E{+}B_1$) & majority cannot couple in \\
\bottomrule
\end{tabular}
\end{table}

\noindent The implication for the whole pipeline is sharp: a target must first be
\emph{admissible} (high $\Aone$ fraction) before any amount of optimisation is
worthwhile. Using the reachability verdict as a front-end filter---rather than
discovering the ceiling empirically after the fact---is, in retrospect, the single
most important methodological lesson of the project.

%------------------------------------------------------------------ 5
\section{Methodology}\label{sec:method}

\subsection{Differentiable orthotropic-plate surrogate}\label{sec:surrogate}
At the heart of the pipeline is a differentiable Kirchhoff plate surrogate
implemented so that the loss is end-to-end differentiable with respect to the
thickness field. We assemble the orthotropic bending-stiffness operator from the
full laminate energy functional~\cite{reddy2003}; each cell may, in the general
formulation, rotate its bending-stiffness tensor by a local angle $\theta$. The
discrete eigenproblem~\eqref{eq:eig} is solved with a sparse solver, and gradients
of the loss with respect to $\matH$ (and $\bm\theta$, $\omega$) are obtained by a
combination of automatic differentiation~\cite{baydin2018,paszke2019} and analytic
eigenpair sensitivities~\cite{fox1968,nelson1976}, the latter being the classical
route to differentiating eigenvalues and eigenvectors of~\eqref{eq:eig}. The
adjoint viewpoint~\cite{giles2000} makes the gradient computation independent of
the number of design variables.

The decisive design decision is one of \emph{role}, not of formulation: the
surrogate is treated as a \textbf{compass, not a judge}. It is fast enough to push
the thickness field and frequency to a ``surrogate-good'' region in seconds, but
it is never allowed to decide the winner; that authority belongs to the
finite-element model. Two opposite misuses---treating the surrogate as the truth,
and treating it as the final referee---were both committed earlier in the project
and both produced over-optimistic designs that did not survive high-fidelity
re-evaluation (Section~\ref{sec:gap}).

\subsection{The anisotropy hypothesis and its collapse}\label{sec:aniso}
A tempting hypothesis motivated the orthotropic formulation: since 3D-printed
parts are naturally anisotropic, letting the fibre orientation $\theta$ vary per
cell would make the bending-stiffness operator cease to commute with $\Dfour$,
allowing a central drive to leak into non-$\Aone$ irreps and thereby break the
ceiling of Section~\ref{sec:theory}. The argument is physically appealing.

A scaling analysis nonetheless predicts that the effect is negligible in the
working band. Writing the forced operator as $\matK-\omega^2\matM$, in the
relevant frequency range the inertial term $\omega^2\lVert\matM\rVert$ exceeds the
stiffness term $\lVert\matK\rVert$ by several orders of magnitude. Because
$\theta$ enters the physics \emph{only} through $\matK$, its influence on the
response is suppressed by the same ratio: the response is, to leading order,
mass-controlled, and the symmetry-breaking that $\theta$ would introduce through
the stiffness is drowned out. The orientation field is, in addition, a
\emph{parasitic} optimisation dimension: it inflates a non-convex loss landscape
and diverts gradient budget away from a clean convergence in $\matH$ and $\omega$.
Independent confirmation that the realised anisotropy ratio of representative
printable materials is far smaller than the optimistic literature values only
reinforces the conclusion. We therefore retain the orthotropic machinery for
generality but treat $\theta$ as effectively frozen, and the isotropic path as the
default. The corresponding quantitative ablation is reported in
Section~\ref{sec:experiments}.

\subsection{A recognisability metric grounded in powder physics}\label{sec:metric}
Early scoring used pixel-overlap measures (IoU/Jaccard, Dice, Chamfer
distance)~\cite{jaccard1912,dice1945,barrow1977}, which are discrete,
non-differentiable, and unstable under mode switching. We replace them with a
metric grounded in the actual deposition physics: powder accumulates where
$|\bm{u}|\approx0$, so the relevant quantity is a powder density concentrated on
nodal lines. The recognisability functional combines (i)~\emph{enrichment}, the
ratio of powder density inside the target region to the global mean;
(ii)~\emph{recall}, the fraction of the target covered; (iii)~\emph{contrast}; and
(iv)~a \emph{direction} term. A continuous, differentiable surrogate of this
functional---an amplitude-valley loss---is used inside the optimiser, while the
discrete powder model is used for evaluation. We also note the connection to
optimal-transport objectives: matching nodal-line measures rather than pixels is
naturally expressed through Sinkhorn divergences~\cite{cuturi2013}, which, like
Wasserstein objectives for waveform matching~\cite{engquist2014}, avoid the
cycle-skipping pathologies of pixelwise $L^2$/IoU losses. The choice of metric is
not cosmetic: it re-shuffles which candidates are considered ``best'' and
therefore steers the entire search, a point we return to in
Section~\ref{sec:discussion}.

\subsection{Two-phase finite-element production pipeline}\label{sec:pipeline}
The production pipeline couples the surrogate to the finite-element model in two
phases that decouple the two things that can go wrong---picking the wrong
frequency and picking the wrong shape.

\paragraph{Phase 1 — modal selection.}
We compute the finite-element eigenfrequencies of the candidate plate, rank each
eigenmode by its likeness to the target, drive the forced
response~\eqref{eq:forced} \emph{on resonance}, and compose the best single mode
or a composite of complementary modes. Driving on resonance is essential: an
off-resonance drive lets the central point force dominate and skews the modes, so
that clean nodal lines never form. We also find, in principle, that an
\emph{interference of two complementary modes} (a beat composite) can be more
target-like than any single eigenmode---a degree of freedom invisible to a
single-mode view.

\paragraph{Phase 2 — trust-region refinement.}
Taking the finite-element response as ground truth, the thickness field is refined
by a trust-region iteration~\cite{conn2000} embedded in a surrogate-management
framework~\cite{booker1999}: the surrogate proposes a step, the finite-element
model re-evaluates it, and the step is accepted or the trust radius shrunk. Modal
correspondence between iterations is tracked by the Modal Assurance
Criterion~\cite{allemang1982}. This explicitly manages the surrogate--FEA gap
through ``trust'' and refines monotonically, in contrast to a naive loop that
corrects frequency and shape simultaneously and diverges.

\subsection{The surrogate--FEA gap}\label{sec:gap}
The thin-plate surrogate and the thick-plate finite-element model differ
permanently: plate theory (thin versus shear-deformable thick), mesh density,
boundary geometry (an abstract central constraint versus a circular clamp), and
two-dimensional plane stress versus a full three-dimensional elasticity tensor.
These are \emph{systematic physical differences}, not noise. Qualitatively, the
surrogate is excellent at low modal order and degrades at high order, and the
mapping between surrogate and finite-element frequencies follows a clean power law
rather than a random scatter---evidence that the gap is structural. The correct
response to a systematic gap is neither to trust the surrogate nor to over-fit
against it, but to use it as a compass and let the finite-element model arbitrate;
this is precisely the rationale for the two-phase pipeline of
Section~\ref{sec:pipeline}. A separate, sobering observation concerns numerical
hygiene: a single cross-scale additive constant in a normalisation
($\text{amp}/(\max(\text{amp})+\varepsilon)$ with a fixed $\varepsilon$) silently
clamped a key metric whenever the finite-element amplitudes were many orders of
magnitude smaller than the surrogate's, systematically \emph{under}-reporting
historical results until corrected.

\subsection{Manufacturable export and software system}\label{sec:software}
The pipeline terminates in a manufacturable artefact. A zero-dependency exporter
converts the $15\times15$ thickness field into a stepped-box STL (one box per
cell, flat-bottomed for bed adhesion, with optional smoothing of the steps),
directly slicable for fused-deposition printing without any orientation control.
The system is delivered as a full-stack application: a multithreaded backend that
orchestrates the surrogate and the finite-element solver (managing the solver
lifecycle, auto-discovering tool paths, and exposing result and diagnostic
endpoints including the reachability verdict of Section~\ref{sec:theory}), and a
single-page frontend organised around target design, material tuning, run control,
results inspection, a uniform-plate pattern gallery filtered by central-drive
excitability, and a 3D-print export tab. The engineering discipline of the project
included a read-only algorithmic audit that quantified cognitive biases in the
metrics, strict roll-back discipline against regressions, and per-experiment
reproducibility.

%------------------------------------------------------------------ 6 (EVOLUTION)
\section{Algorithmic evolution: three eras and roughly twenty lines}\label{sec:eras}

The architecture of Section~\ref{sec:method} did not arrive by design; it
precipitated out of an extended trial-and-error process that is itself part of the
contribution, because each abandoned line marks a wall of the feasible region. We
organise the history into three eras followed by a modern production period
(Table~\ref{tab:eras}). The milestone numbers quoted below describe the
\emph{development trajectory}; the controlled experimental validation remains the
(deliberately empty) subject of Section~\ref{sec:experiments}.

\begin{table}[h]
\centering
\caption{The three eras and the modern production period.}
\label{tab:eras}
\small
\begin{tabular}{@{}p{2.9cm}p{3.0cm}p{5.2cm}p{3.4cm}@{}}
\toprule
Era & Phases & Core paradigm & Outcome \\
\midrule
I --- Blind search & Phase A--I (9 lines) & ``guess thickness $\to$ run FEA $\to$ tweak'' & score stuck at $0.034$; all below bar \\
II --- Structured & W1--W6 (6 lines) & prior filter + main-loss redesign + gradient path & IC score pushed to $0.171$ \\
III --- Recognisability & W7--W8 (2 lines) & redefine success via powder physics & X-shape enrichment $3.76\times$ (ceiling) \\
Modern production & W9, W10, two-phase ($\sim$3 lines) & differentiable surrogate + FEA trust-region loop + productisation & two-phase pipeline + software \\
\bottomrule
\end{tabular}
\end{table}

\subsection{Era I --- blind search (Phases A--I)}\label{sec:eraI}
The most naive formulation treats the $225$-dimensional thickness field directly
as the optimisation variable and searches it with gradient-free methods
(Table~\ref{tab:eraI}). Random and evolutionary search only ever produced the
``generic plate-mode family'' (arcs, stars, rings, centre-dominant blobs) with no
letter-like structure; response-guided error maps could say \emph{where} a pattern
was wrong but not \emph{which} nodal line to move where, because eigenmodes are
global objects; a first Kirchhoff--Love proxy already exposed the
surrogate-to-FEA gap; and a strict-scoring upgrade honestly cut an inflated best
of $\approx0.2$ down to $0.034$. Expanding the design contract from $225$ to
$675$ dimensions did not help. Era~I ended with a best honest score of $0.034$,
far below the $0.08$ engineering bar, and a clear diagnosis: more generations
would not help; the problem had to be redesigned from the physics and mathematics
upward.

\begin{table}[h]
\centering
\caption{Era I --- the nine blind-search lines and their lessons.}
\label{tab:eraI}
\small
\begin{tabular}{@{}p{0.7cm}p{5.3cm}p{2.2cm}p{6.8cm}@{}}
\toprule
Line & Method & Methodological basis & Lesson \\
\midrule
A & Random / evolutionary search (GA, mutation, neighbour-diff repair, skeleton guidance) & \cite{holland1975,bendsoe2003} & Only the generic plate-mode family; no letters \\
B & Response-guided loop (missing/extra error maps; iterative inversion) & \cite{tarantola2005} & Error maps locate error but not the corrective move; modes are global \\
C & Kirchhoff--Love plate proxy (sparse biharmonic eigensolve) & \cite{leissa1969,waller1939} & First exposure of the surrogate-to-FEA gap; proxy scores do not transfer \\
D & Strict shape scoring ($v1\to v5$; IoU/Dice/Chamfer) & \cite{jaccard1912,dice1945,barrow1977} & Honest metrics beat pretty ones: best falls from $\approx0.2$ to $0.034$ \\
E & Design-contract expansion (SIMP density/loss/material, $225\to675$) & \cite{bendsoe1989,bendsoe1999} & More dimensions are not the answer; placement still fails \\
F & Auxiliary-physics local heuristics (direct/pattern search) & \cite{hooke1961} & Most robust historical source, but capped at $\approx0.032$--$0.034$ \\
G & Latent joint physics (hand-designed basis; POD) & \cite{berkooz1993} & A hand-made basis can miss key directions \\
H & CMA-ES in latent coefficient space & \cite{hansen2001} & Exposes surrogate optimism: good on proxy, poor on FEA \\
I & Response-calibrated scheduling (surrogate/Bayesian optimisation) & \cite{jones1998,shahriari2016} & First systemic treatment of ``surrogate is unreliable''; not fully landed \\
\bottomrule
\end{tabular}
\end{table}

\subsection{Era II --- structured search (W1--W6)}\label{sec:eraII}
The second era replaced blind search with structure: a physical prior filter, a
differentiable main loss, a connected gradient path, and a data-driven design
subspace (Table~\ref{tab:eraII}). Two points deserve emphasis. First, the
realizability prior of W1 already \emph{predicted the entire IC outcome} through a
$\Dfour$-mismatch of $0.78$; that this verdict was not used as a front-end filter
at the time is, in retrospect, the single biggest methodological miss of the
project (it is precisely the theory later formalised in
Section~\ref{sec:theory}). Second, W6 --- the climax of the era --- pushed the IC
score from $0.034$ to $0.171$, a roughly fivefold gain that nonetheless still did
not look like IC, which is exactly what motivated the redefinition of success in
Era~III.

\begin{table}[h]
\centering
\caption{Era II --- the six structured lines and their key results.}
\label{tab:eraII}
\small
\begin{tabular}{@{}p{0.7cm}p{5.3cm}p{2.2cm}p{6.8cm}@{}}
\toprule
Line & Method & Methodological basis & Key result / lesson \\
\midrule
W1 & Realizability prior filter (Courant nodal domains + $\Dfour$ mismatch + stroke width) & \cite{courant1953,pleijel1956,waller1939} & IC mismatch $0.78$ predicted the ending; not used up front (biggest miss) \\
W2 & Main-loss redesign (IoU $\to$ amplitude valley) & \cite{nocedal2006} & Smoother optimisation; physical ceiling unchanged \\
W3 & Differentiable plate + adjoint ($\partial\,\text{loss}/\partial\matH$) & \cite{giles2000,baydin2018,paszke2019,fox1968,nelson1976} & Gradient path finally connected; first structural upgrade \\
W4 & Low-dimensional manifold search (PCA of candidates) & \cite{pearson1901,berkooz1993} & Data basis beats hand-made basis; markedly lower loss than W3 \\
W5 & Homotopy continuation (warm-start toward target) & \cite{allgower1990,blake1987} & Stalls at $\lambda\approx0.3$ --- an honest reachability indicator \\
W6 & Spectral placement (cluster eigenvalues near drive) & \cite{pedersen2000,achtziger2007} & Era II climax: IC score $0.034\to0.171$ ($5\times$), still not legible \\
\bottomrule
\end{tabular}
\end{table}

\subsection{Era III --- recognisability (W7--W8)}\label{sec:eraIII}
Era~III was triggered by the realisation that a good score ($0.171$) still did not
look like the target. Success was redefined to accept ``similar, not identical'',
and the scoring system was rewritten around powder physics. The multi-frequency
root-mean-square composition of W7, optimised with a first-order stochastic
method~\cite{kingma2015} over the forced response~\cite{leissa1969}, proved a
\emph{pseudo}-breakthrough: the optimiser learns to weight $|\bm{u}|^2$ to
fabricate low-amplitude regions that the finite-element model cannot reproduce,
and it tends to collapse back to a single frequency. W8 then replaced pixel
overlap with the powder-deposition metric of Section~\ref{sec:metric}, formulated
through an optimal-transport view of nodal-line matching~%
\cite{chladni1787,cuturi2013,engquist2014}. Re-ranking the entire candidate history under the new
metric revealed that the project had optimised the \emph{wrong} metric for months
and that better candidates had been buried. The X-shape enrichment reached
$2.84\times$, and $3.76\times$ after the normalisation fix discussed below, which
became the project ceiling.

\subsection{Modern production --- W9, W10 and the two-phase pipeline}\label{sec:modern}
The modern period built the pieces that the earlier eras had only gestured at. The
W9 trust-region loop was structurally correct but under-tuned (too few inner
steps, too-weak correction, a trust radius that shrank to $0.094\,\text{mm}$), and
its best finite-element result of $1.40\times$ did not break through --- the lesson
being that circling a poor surrogate does not help when the root cause is the
surrogate physics itself; W9 nonetheless contributed the trust-region and
surrogate-management machinery with modal tracking~%
\cite{conn2000,booker1999,allemang1982} that Phase~2 later inherited. W10
implemented the differentiable orthotropic plate of
Section~\ref{sec:surrogate}~\cite{reddy2003}, and the two-phase pipeline of
Section~\ref{sec:pipeline} finally realised the iterative calibration loop that
Era~I had owed all along, by decoupling the frequency correction (driven directly
at the finite-element resonance) from the shape correction (a small trust-region
residual).

\subsection{Cross-cutting empirical findings}\label{sec:crosscut}
Three findings recur across the eras and deserve separate statement.

\paragraph{On-resonance versus off-resonance drive.}
A late, bug-level discovery: Phase~1 once defaulted to driving at
$f_{\text{eig}}+1.5\,\text{Hz}$ (off-resonance) to avoid a solver singularity, but
this let the central point force dominate the response and skewed the modes, so
that cross- and X-patterns never formed. Switching to an on-resonance drive
(offset zero) instantly recovered clean modal nodal lines. The off-resonance
artefact also explains a spurious ``bar through the middle'': off resonance, the
centre is an antinode bright streak pushed up by the point force, not a true nodal
line.

\paragraph{A cross-scale numerical constant once hid the truth.}
The powder-density normalisation $\text{amp}/(\max(\text{amp})+\varepsilon)$ with a
fixed $\varepsilon$ was harmless in surrogate-space, where amplitudes are
$\mathcal{O}(1)$, but in finite-element space, where amplitudes are
$\mathcal{O}(10^{-12})$, the constant was a thousand times larger than the signal
and silently clamped the key metric to $\approx 1.00$. After the fix, all
historical finite-element numbers rose (the X-shape from $2.84\times$ to
$3.76\times$). The lesson is blunt: a cross-scale additive constant is a silent
killer.

\paragraph{The metric defines the direction of search.}
Five primary metrics were used over the project, and each switch re-shuffled the
``best candidate'' (one W6 variant was first under one metric and tenth under
another): pixel overlaps (IoU/Dice/Chamfer); a smoothed pixel score; an
amplitude-valley loss; a multi-frequency RMS (easy for the optimiser to cheat);
and finally the powder-physics recognisability of Section~\ref{sec:metric}. A
separate read-only algorithmic audit later quantified residual biases even in the
last family (for example, that enrichment can prefer a compact central blob over an
extended shape, and that fixed-percentile thresholds are sensitive to powder
density). The meta-lesson, which we restate in Section~\ref{sec:discussion}, is to
define ``success'' before optimising it.

%------------------------------------------------------------------ 7  (BLANK)
\section{Experimental setup and results}\label{sec:experiments}

\noindent\emph{This section is intentionally left as a scaffold; the quantitative
experimental write-up is pending and will populate the subsections below. No
numerical results are reported here yet.}

\subsection{Experimental setup}
\noindent\textit{[To be completed.]} Plate and material parameters, the
finite-element model configuration (mesh, element type, damping, eigenfrequency
and frequency-domain solver settings via LiveLink), the powder model parameters,
optimiser hyper-parameters, and the target set.

\subsection{Validation of the reachability theory}
\noindent\textit{[To be completed.]} Empirical confirmation that achievable
recognisability tracks the analytic $\Aone$ fraction of Table~\ref{tab:reach}
across the target set.

\subsection{Anisotropy ablation (material battery)}
\noindent\textit{[To be completed.]} A target $\times$ material ablation comparing
an anisotropic configuration (orientation active) against an isotropic baseline
(orientation frozen), together with the direct sensitivity measurement of the
composite amplitude to $\theta$.

\subsection{Quantification of the surrogate--FEA gap}
\noindent\textit{[To be completed.]} Modal-assurance statistics across modal
order, the surrogate-to-finite-element frequency map, and the fraction of the
surrogate-predicted improvement that the finite-element pipeline reproduces.

\subsection{End-to-end case studies}
\noindent\textit{[To be completed.]} Full pipeline results for representative
targets (X-shape, IC, cross, single diagonal, ring), including the best
single-mode and composite drives, the refined thickness field, and the exported
printable geometry.

%------------------------------------------------------------------ 7
\section{Discussion}\label{sec:discussion}

The four findings of the abstract cohere into a single design philosophy. First,
the reachability theory of Section~\ref{sec:theory} reframes failure: when a
target such as IC cannot be produced, this is not an algorithmic shortfall but a
structural fact about which irreps a central drive can excite, and it is better to
state that fact precisely than to keep optimising against an inadmissible target.
Second, the anisotropy analysis of Section~\ref{sec:aniso} is a cautionary tale
about ``physically plausible'' degrees of freedom: a mechanism can be real
(orientation does break symmetry through the stiffness) and yet be quantitatively
irrelevant (because the response is mass-controlled in the working band), so a
favourite hypothesis must be allowed to die when the scaling says so. Third, the
surrogate--FEA gap of Section~\ref{sec:gap} shows that the operative ceiling in
practice is often not the deep physics but the fidelity gap between the model one
optimises and the model one believes, and that the disciplined response is to
assign each model the role it is good at. Fourth, the metric discussion of
Section~\ref{sec:metric} underlines that in an inverse problem the definition of
``success'' is itself a design choice that silently steers the search; defining it
correctly---here, in terms of powder physics rather than pixel overlap---is
prerequisite to optimising at all. The recurring meta-lesson is that a substantial
fraction of effort in a fixed-hardware inverse problem is well spent on
\emph{characterising the boundary} rather than on pushing against it.

%------------------------------------------------------------------ 8
\section{Limitations and future work}\label{sec:limits}

The limitations are inherent rather than incidental. Any target that is not
$\Dfour$-compatible---an arbitrary letter pair, a single diagonal, an L-shape, a
generic logo---is physically unreachable under a central point drive, and this is a
hard ceiling rather than a tuning problem. Material anisotropy offers little
headroom, both because realised anisotropy ratios are modest and because the
orientation gradient is negligible off resonance. A surrogate is never the
finite-element model, so any surrogate-only ``breakthrough'' must be re-checked at
high fidelity.

It follows that genuine further progress requires changing the
\emph{hardware}, in decreasing order of leverage: a $\Dfour$-symmetry-breaking
excitation (multiple shakers or an off-centre drive), a changed boundary (a
non-square plate), or discrete topology (holes and slots). These lie outside the
fixed contract of the present study but are the only routes that can lift the
$\Aone$ ceiling. Within the present contract, promising refinements include a
symmetry-aware direction metric (to fix principal-axis degeneracy on symmetric
targets), a genuinely physical time-division multi-frequency drive rather than a
numerical root-mean-square composite, and further tuning of the trust-region
parameters of Phase~2.

%------------------------------------------------------------------ 9
\section{Conclusion}\label{sec:conclusion}

Under a fixed hardware contract---a square plate with a single central
excitation---we have shown that the right question is not ``how strong can the
inversion be'' but ``where is the reachability boundary, and why''. A
symmetry-based theory identifies the $\Aone$ irreducible representation as the only
channel a central drive can excite and turns this into a closed-form, target-specific
ceiling; a scaling analysis demotes material anisotropy from a hoped-for escape
route to a parasitic degree of freedom; and a two-phase, surrogate-as-compass
pipeline manages the dominant practical obstacle, the gap between a thin-plate
surrogate and a thick-plate finite-element model. The deliverable is an honest,
reproducible, manufacturable system rather than an impossible promise of arbitrary
inversion. The remaining headroom, we argue, is not in the algorithm but in the
hardware. The experimental validation that quantifies these claims is reported
separately; the corresponding section here is left as a scaffold by design.

%------------------------------------------------------------------ Ack
\section*{Acknowledgements}
This work was carried out as a first-year undergraduate Term~3 summer project at
Imperial College London. We thank the maintainers of the open-source scientific
software on which the pipeline is built.

%------------------------------------------------------------------ References
\begin{thebibliography}{99}
\setlength{\itemsep}{2pt}

\bibitem{chladni1787} E.~F.~F. Chladni. \textit{Entdeckungen \"uber die Theorie
des Klanges}. Leipzig, 1787.

\bibitem{waller1939} M.~D. Waller. Vibrations of free square plates.
\textit{Proc. Phys. Soc.}, 51:831, 1939.

\bibitem{reissner1945} E. Reissner. The effect of transverse shear deformation on
the bending of elastic plates. \textit{J. Appl. Mech.}, 12:A69--A77, 1945.

\bibitem{mindlin1951} R.~D. Mindlin. Influence of rotatory inertia and shear on
flexural motions of isotropic, elastic plates. \textit{J. Appl. Mech.},
18:31--38, 1951.

\bibitem{leissa1969} A.~W. Leissa. \textit{Vibration of Plates}. NASA SP-160,
1969.

\bibitem{reddy2003} J.~N. Reddy. \textit{Mechanics of Laminated Composite Plates
and Shells: Theory and Analysis}, 2nd ed. CRC Press, 2003.

\bibitem{tinkham1964} M. Tinkham. \textit{Group Theory and Quantum Mechanics}.
McGraw--Hill, 1964.

\bibitem{courant1953} R. Courant and D. Hilbert. \textit{Methods of Mathematical
Physics}, Vol.~1. Interscience, 1953.

\bibitem{pleijel1956} \r{A}. Pleijel. Remarks on Courant's nodal line theorem.
\textit{Comm. Pure Appl. Math.}, 9:543--550, 1956.

\bibitem{letcher2014} T. Letcher and M. Waytashek. Material property testing of
3D-printed specimen in PLA on an entry-level 3D printer. \textit{ASME IMECE},
2014.

\bibitem{pyl2018} L. Pyl, K.-A. Kalteremidou, and D. Van Hemelrijck. Exploration
of the mechanical properties of FDM-printed PLA. \textit{Procedia Manufacturing},
14:104--, 2018.

\bibitem{formlabs2018} Formlabs. \textit{Validating Isotropy in SLA 3D Printing}.
Technical white paper, 2018.

\bibitem{vatphoto2025} Vat-photopolymerisation anisotropy study.
\textit{Scientific Reports}, 15:97294, 2025.

\bibitem{bambutds} Bambu Lab. \textit{PETG-CF Technical Data Sheet} (TDS V3.0).

\bibitem{holland1975} J.~H. Holland. \textit{Adaptation in Natural and Artificial
Systems}. University of Michigan Press, 1975.

\bibitem{bendsoe2003} M.~P. Bends\o{}e and O. Sigmund. \textit{Topology
Optimization: Theory, Methods and Applications}. Springer, 2003.

\bibitem{tarantola2005} A. Tarantola. \textit{Inverse Problem Theory and Methods
for Model Parameter Estimation}. SIAM, 2005.

\bibitem{jaccard1912} P. Jaccard. The distribution of the flora in the alpine
zone. \textit{New Phytologist}, 11:37--50, 1912.

\bibitem{dice1945} L.~R. Dice. Measures of the amount of ecologic association
between species. \textit{Ecology}, 26:297--302, 1945.

\bibitem{barrow1977} H.~G. Barrow et al. Parametric correspondence and chamfer
matching. \textit{Proc. IJCAI}, 1977.

\bibitem{bendsoe1989} M.~P. Bends\o{}e. Optimal shape design as a material
distribution problem. \textit{Struct. Optim.}, 1:193--202, 1989.

\bibitem{bendsoe1999} M.~P. Bends\o{}e and O. Sigmund. Material interpolation
schemes in topology optimization. \textit{Arch. Appl. Mech.}, 69:635--654, 1999.

\bibitem{hooke1961} R. Hooke and T.~A. Jeeves. ``Direct search'' solution of
numerical and statistical problems. \textit{J. ACM}, 8:212--229, 1961.

\bibitem{berkooz1993} G. Berkooz, P. Holmes, and J.~L. Lumley. The proper
orthogonal decomposition in the analysis of turbulent flows. \textit{Annu. Rev.
Fluid Mech.}, 25:539--575, 1993.

\bibitem{hansen2001} N. Hansen and A. Ostermeier. Completely derandomized
self-adaptation in evolution strategies. \textit{Evol. Comput.}, 9(2):159--195,
2001.

\bibitem{jones1998} D.~R. Jones, M. Schonlau, and W.~J. Welch. Efficient global
optimization of expensive black-box functions. \textit{J. Global Optim.},
13:455--492, 1998.

\bibitem{shahriari2016} B. Shahriari et al. Taking the human out of the loop: A
review of Bayesian optimization. \textit{Proc. IEEE}, 104:148--175, 2016.

\bibitem{nocedal2006} J. Nocedal and S.~J. Wright. \textit{Numerical
Optimization}, 2nd ed. Springer, 2006.

\bibitem{giles2000} M.~B. Giles and N.~A. Pierce. An introduction to the adjoint
approach to design. \textit{Flow Turbul. Combust.}, 65:393--415, 2000.

\bibitem{baydin2018} A.~G. Baydin et al. Automatic differentiation in machine
learning: a survey. \textit{J. Mach. Learn. Res.}, 18:1--43, 2018.

\bibitem{paszke2019} A. Paszke et al. PyTorch: An imperative style, high-performance
deep learning library. \textit{Proc. NeurIPS}, 2019.

\bibitem{fox1968} R.~L. Fox and M.~P. Kapoor. Rates of change of eigenvalues and
eigenvectors. \textit{AIAA J.}, 6:2426--2429, 1968.

\bibitem{nelson1976} R.~B. Nelson. Simplified calculation of eigenvector
derivatives. \textit{AIAA J.}, 14:1201--1205, 1976.

\bibitem{pearson1901} K. Pearson. On lines and planes of closest fit to systems of
points in space. \textit{Phil. Mag.}, 2:559--572, 1901.

\bibitem{allgower1990} E.~L. Allgower and K. Georg. \textit{Numerical Continuation
Methods: An Introduction}. Springer, 1990.

\bibitem{blake1987} A. Blake and A. Zisserman. \textit{Visual Reconstruction}.
MIT Press, 1987.

\bibitem{pedersen2000} N.~L. Pedersen. Maximization of eigenvalues using topology
optimization. \textit{Struct. Multidisc. Optim.}, 20:2--11, 2000.

\bibitem{achtziger2007} W. Achtziger and M. Ko\v{c}vara. On the maximization of
the fundamental eigenvalue in topology optimization. \textit{Struct. Multidisc.
Optim.}, 34:181--195, 2007.

\bibitem{kingma2015} D.~P. Kingma and J. Ba. Adam: A method for stochastic
optimization. \textit{Proc. ICLR}, 2015.

\bibitem{cuturi2013} M. Cuturi. Sinkhorn distances: Lightspeed computation of
optimal transport. \textit{Proc. NeurIPS}, 2013.

\bibitem{engquist2014} B. Engquist and B.~D. Froese. Application of the
Wasserstein metric to seismic signals. \textit{Commun. Math. Sci.}, 12:979--988,
2014.

\bibitem{conn2000} A.~R. Conn, N.~I.~M. Gould, and P.~L. Toint.
\textit{Trust-Region Methods}. SIAM, 2000.

\bibitem{booker1999} A.~J. Booker et al. A rigorous framework for optimization of
expensive functions by surrogates. \textit{Struct. Multidisc. Optim.}, 17:1--13,
1999.

\bibitem{allemang1982} R.~J. Allemang and D.~L. Brown. A correlation coefficient
for modal vector analysis. \textit{Proc. Int. Modal Analysis Conf. (IMAC)}, 1982.

\bibitem{acta2024} Topology optimization design of Chladni patterns for vibration
mode manipulability. \textit{Acta Mechanica Sinica}, 2024.
DOI:~10.1007/s10409-023-23445-x.

\bibitem{tcherniak2002} D. Tcherniak. Topology optimization of resonating
structures using SIMP method. \textit{Int. J. Numer. Meth. Eng.}, 54:1605--1622,
2002.

\bibitem{bellette} P. Bellette. \texttt{chladni\_inverse\_design}. Open-source
implementation (energy-landscape formulation).

\bibitem{misseroni2016} D. Misseroni, A.~B. Movchan, and N.~V. Movchan. Cymatics
for the cloaking of flexural vibrations in a structured plate. \textit{Scientific
Reports}, 6:23929, 2016.

\bibitem{tuanchen} P.-T. Tuan, Y.-F. Chen, et al. Maximum-entropy reconstruction
of Chladni forced-response nodal patterns via the inhomogeneous Helmholtz
equation.

\bibitem{zhou2016} Q. Zhou et al. Controlling the motion of multiple objects on a
Chladni plate. \textit{Nature Communications}, 7:12764, 2016.

\bibitem{allaire2004} G. Allaire, F. Jouve, and A.-M. Toader. Structural
optimization using sensitivity analysis and a level-set method. \textit{J.
Comput. Phys.}, 194:363--393, 2004.

\end{thebibliography}

\end{document}
